\section{Seven Gifts of the Holy Spirit}

\begin{quotex}
Esoteric science in its teaching goes far beyond simple information. Its purpose, in fact, is nothing less than the transformation of the very being of those who study it \flright{\textsc{Boris Mouravieff}}

\end{quotex}
We can represent the Fallen World, i.e., the world in which mankind determines what is right and wrong, true and false, as a closed vector space, in which each vector represents an idea or meme. Depending then on your specific position in the space, you will be subjected to various idea vectors of different strengths and directions, randomly with no fixed pattern. Hence, there is no true standard of right and wrong, just various forces acting on your mind. Most of the inhabitants of this world accept it as it is, while convincing themselves that their particular vectors are somehow truer than other vectors. We can refer to these as the A influences.

Since that can never be more than a war of all against all, we can ask if the world process is truly closed.

If the world process is not closed, then higher influences can enter into it. Specifically, there are the B influences that come from valid (or formerly valid) traditions. Unlike the A influences, the B influences all have the same direction. Furthermore, there are even higher influences such as esoteric teachings and even angelic (or planetary) influences. It is possible to become attuned to such higher influences in order to escape the influences of the world process. The highest of such influences will come from the Holy Spirit.

The Reformation denied such higher influences, which are expressed in symbolic form, and asserted that only the literally physical and historical interpretations of scriptures are valid. They also denied the authority of philosophy and metaphysics. This worldview is not something to debate; the real problem is that at that time after the Medieval period, fewer men were able to even understand the more ancient esoteric and symbolic understanding of religion. That is, they could perhaps grasp the B influences, but not the C or higher influences. Ultimately, physics and history, profane sciences which deal solely with the A influences, became the sole source of authority. The esoteric sciences were forgotten or driven underground.

\paragraph{The Seven Gifts}
The seven gifts of the Holy Spirit are given (or at least strengthened) in Confirmation. Unfortunately, in most cases these gifts are ineffectual, and are hardly known. Hence they remain virtual, not fully actualized. In addition, over time the meaning of words often change. So it is worth the trouble to understand how certain concepts were used in Medieval Christendom. English is inadequate as a language for metaphysics, so we will sometimes resort to the Latin or Greek equivalents for clarity.

First of all, these are the gifts of the Holy Spirit:

\begin{itemize}
\item Wisdom (Sophia) 
\item Intelligence (intellectus, nous) 
\item Counsel 
\item Strength 
\item Knowledge (gnosis) 
\item Piety 
\item Fear of God 
\end{itemize}
These gifts require a period of purification in order to become fully active. Few today are willing to, and often do not even how to, engage in that process of purification.

\paragraph{Wisdom}
We have addressed the topic of Wisdom, Sophia, Chokmah, have been often addressed. For a good description of Sophia, especially how it was described by \textbf{Vladimir Solovyov}, please visit \textit{The Immaculate Conception and the Divine Sophia}\footnote{\url{https://crc-internet.org/our-doctrine/national-restoration/holy-russia/the-immaculate-conception-and-the-divine-sophia.html?r=soloviev}}.

\paragraph{Intelligence}
According to St. Anthony, only the rational man is in union with God (or Nous in Platonic or Aristotelian terms), and reason is not at all common to all men. The rational faculty unites man with divine power, power over oneself, one's thoughts, one's acts. As such, it is a faculty that must be developed\footnote{\textit{Rational Man and Moral Law}, see Section \ref{sec:RationalMalMoralLaw}.}.

Man is dominated by three psychic forces: eros (desire, concupiscence), thymos (spiritedness, emotion), and nous (intellect). Most people determine their goals based on desire and emotion. The intellect merely serves an instrumental function, i.e., it determines the means to reach the goals defined by eros and thymos.

However, real intelligence needs to determine man's proper ends or goals, not just the means. This gift of the Holy Spirit enables us to choose our proper goals based on intelligence, rather than on concupiscence or emotion.

\paragraph{Counsel}
Natural willing designates the movement of a creature in accordance with the principle (\emph{logos}) of its nature towards the fulfilment (\emph{telos}, \emph{stasis}) of its being. (Wikipedia)

Gnomic will, on the other hand, requires a process of deliberation and decision making. It is therefore dualistic, dual in the evil sense described by \textbf{Valentin Tomberg}. If purity of heart is to will one thing, then our very being will spontaneously (NOT impulsively) choose its proper end. Specifically, if we rely on the Counsel of the Holy Spirit, rather than on our own Will, we will gain the power to fulfill our purposes in life.

\paragraph{Strength}
The concept of Strength is that of an intermediary between pure consciousness and manifestation. It is the link between the idea and the phenomenon\footnote{Letter XI from \emph{Meditations on the Tarot}.}.

The rational man is not just a learned man or scholar, but rather someone who not only knows good from bad, but also has the power to act on it. There are three faults that beset man from the Fall: concupiscence, ignorance, malice. This is not an abstract dogma, but something that is readily observable, both in our own selves and certainly in other people. As we pointed out in \textit{Space Aliens, Pigs, Roosters, and Snakes}\footnote{\url{https://www.gornahoor.net/?p=9647}}, even the Buddhists became aware of these faults, although they did not know exactly how they arose.

One of the obstacles of the spiritual life is that even when we become aware of these faults within us, we still find it difficult to overcome them. Weakness, then, is the fourth fault which prevents us from doing what we need to do.

Hence, the Holy Spirit gives us the Strength to overcome the faults arising from the Fall.

\paragraph{Knowledge}
There are three levels of knowledge: doxa, dianoia, episteme. Doxa is opinion, or knowledge of sensory things. Dianoia refers to logical, hypothetical, speculative, scholarly, and argumentative knowledge, while episteme is direct, intuitive knowledge of spiritual things.

Intelligent people are seldom able to surpass dianoia. They think that everything is open to debate or that the next book will reveal the secret. They often know the writings of Saints and Mystics, yet don't bother to achieve the higher states of those saints.

The knowledge, i.e., gnosis, that results from episteme is of a totally different order. It is direct and certain. Episteme does not have a natural source, but rather is a Gift of the Holy Spirit.

\paragraph{Piety}
The cause of the Fall was the desire to follow one's own desires rather than to be obedient to God's Will. This Gift of the Holy Spirit overcomes the Fall and inspires us to obedience. We learn to love to serve God's will.

We find it easier. And faith becomes stronger. No longer are we confused about dogmas, especially when we notice the strange beliefs of irreligious moderns.

\paragraph{Fear of God}
\begin{quotex}
This moral instrument is moulded to fit the traditions of milieu and family. It is shaped from birth onwards by education. Clearly, without this instrument, the organization of social life in all its forms is unthinkable. Yet because of its nature, it cannot guarantee good and equitable human conduct; to ensure its own existence in times of peace, human society has always been obliged to have recourse to constraint and the application of penalties: necessary remedies, since morality will never be strong enough to curb the extreme and anarchistic tendencies of the Personality. The latter, in effect, lacks that kind of consciousness that the practical studies of religion describe as the fear of God. \flright{\textsc{Boris Mouravieff}}

\end{quotex}
Even if the Fear of God is an imperfect motivation to do God's Will, it is certainly praiseworthy. It leads to repentance and the desire to improve. Most of those stuck in the world process fear the government or social disapproval rather than God; they then seldom seek to improve. That is man's natural condition, therefore this is a great Gift of the Holy Spirit.


\hfill

\emph{But God knows best.}



\flrightit{Posted on 2019-12-06 by Cologero }

\begin{center}* * *\end{center}

\begin{footnotesize}\begin{sffamily}



\texttt{S.P on 2020-08-09 at 21:47 said: }

Since you mention the dogmas, how should we understand them in relation to a Primordial Tradition? If the Church demands the belief that outside of it there is no salvation, at the risk of mortal sin, how to reconcile this with the acceptance of other religions and a Primordial Tradition? Is such a dogma a mere necessity for orthodoxy to protect itself from corruption?

PS: I'm sorry for the possible bad English.


\hfill

\texttt{Cologero on 2020-09-11 at 16:42 said: }

Obviously, there was a Primordial Tradition, known to Noah, who would have passed it on, right up to Abram. When the nations were scattered after Babel, the various religions were created. No one is predestined to hell (also a dogma), therefore, as Augustine asserts, everyone is given what he needs for salvation. (The details are not for us to know.) Moreover, as Augustine asserts, there is but One Tradition, now known as Catholicism.

This ends the discussion.


\hfill

\texttt{S.P. on 2020-09-11 at 20:52 said: }

Well, thanks for the reply, Cologero. I suppose it relates to the conception of a primordial monotheism already discussed earlier.

A few months ago, through reading, I obtained a particular state of inspiration, which I attribute (without going too deep) to the Holy Spirit. It is a very curious experience, there is a feeling of lightness and clarity of mind, and a natural disposition (if the term is allowed) towards the Transcendent. Unfortunately, human weakness is very efficient. Of course, you have what you deserve.


\end{sffamily}\end{footnotesize}
